\section*{الفصل \ref{ch:b3:planetary-atmospheres} --- الغلاف الجوي للكواكب وصلاحيتها للحياة}

\begin{solution}{exo:b3:planetary-atmospheres:1}
$H = RT/Mg$: الزهرة $8.314\times737/(0.04401\times8.87) = \qty{15.7}{km}$؛
الأرض $8.314\times288/(0.02897\times9.80) = \qty{8.4}{km}$؛ المريخ
$8.314\times214/(0.04334\times3.71) = \qty{11.1}{km}$؛ تيتان
$8.314\times94/(0.0280\times1.35) = \qty{20.7}{km}$. الغلاف الجوي لتيتان،
على برودته، هو الأوسع امتدادًا: جاذبيته أضعف بسبع مرات
من جاذبية الأرض.
\end{solution}

\begin{solution}{exo:b3:planetary-atmospheres:2}
$\lambda = GMm/kTr = \num{6.674e-11}\times\num{5.97e24}\times\num{1.674e-27}
/(\num{1.381e-23}\,T\times\num{6.878e6})$: 8.78 عند \qty{800}{K}، 4.39 عند
\qty{1600}{K}. الكسران المتسرّبان هما $9.78\,e^{-8.78} =
\num{1.5e-3}$ و$5.39\,e^{-4.39} = \num{6.7e-2}$: الدورة الشمسية
تغيّر الكسر المتسرّب للهيدروجين بعامل يقارب 44.
\end{solution}

\begin{solution}{exo:b3:planetary-atmospheres:3}
(أ)~$d = S_\mathrm{eff}^{-1/2}$: \omterm{def:b3:planetary-atmospheres:runaway}{الدفيئة الرطبة} \qty{0.993}{au}، \omterm{def:b3:planetary-atmospheres:runaway}{الدفيئة الجامحة}
\qty{0.975}{au}، الدفيئة القصوى \qty{1.705}{au}. (ب)~الزهرة، عند 1.91،
تتلقّى أكثر من دفق \omterm{def:b3:planetary-atmospheres:runaway}{الدفيئة الجامحة}، بل أكثر من الحدّ التجريبي
للزهرة الحديثة (1.78): فهي تقع داخل الحافة الداخلية. المريخ، عند 0.43،
يقع بين 1.014 و0.3438، داخل النطاق الصالح للحياة المحافظ:
جفافُه يعود إلى صغر حجمه وفقدان غلافه الجوي، لا
إلى بُعده.
\end{solution}

\begin{solution}{exo:b3:planetary-atmospheres:4}
(أ)~$\log_{10}p = 3.9895 - 443.028/93.66 = -0.741$، $p =
\qty{0.18}{bar}$. (ب)~$0.05\times1.6 = \qty{0.080}{bar}$، أي رطوبة نسبية
تبلغ $0.080/0.18 = 44\,\%$. (ج)~الهواء فوق البحار بعيد جدًّا عن
التشبّع: يتبخّر الميثان منها، كما يتبخّر الماء من محيطات
الأرض، ويجب أن يرتفع الهواء الرطب ويبرد قبل أن يمطر.
\end{solution}

\begin{solution}{exo:b3:planetary-atmospheres:5}
$\lambda_\mathrm D = 7.02\times2.014/1.008 = 14.0$. $u_\mathrm D/u_\mathrm H
= \sqrt{1.008/2.014} = 0.71$. النسبة لكل ذرّة $= 0.71\times(15.0\,
e^{-14.0})/(8.02\,e^{-7.02}) = 0.71\times\num{1.69e-3} = \num{1.2e-3}$.
يكاد \omterm{def:b3:planetary-atmospheres:jeans-escape}{إفلات جينز} لا يسمح بخروج الديوتيريوم: إنه مرشّح قويّ، مع
$\alpha \approx 0.001$.
\end{solution}

\begin{solution}{exo:b3:planetary-atmospheres:6}
(أ)~$0.016/\num{1.5575e-4} = 103$. (ب)~عند $\alpha = 0$ يكون الإثراء
$M_0/M$: $M_0 = 103\times\num{6e15} = \qty{6.2e17}{kg}$، أو
$\num{6.2e17}/(1000\times\num{4.60e14}) = \qty{1.3}{m}$ من الماء فوق
الزهرة. (ج)~$M_0/M = 103^{1/0.2} = \num{1.2e10}$، أي مخزون سابق من
\qty{7e25}{kg}، أي اثنا عشر ضعف كتلة الأرض: وهذا مستحيل. يحتاج
الإثراء إلى إفلات شديد التجزّؤ (\omterm{def:b3:planetary-atmospheres:jeans-escape}{إفلات جينز} أو \omterm{def:b3:planetary-atmospheres:hydrodynamic}{الإفلات غير الحراري})،
نشطًا على الأقل خلال فقدان آخر ما تبقّى من الماء، أو الحالة
المستقرّة مع إمداد بالماء.
\end{solution}

\begin{solution}{exo:b3:planetary-atmospheres:7}
(أ)~في أربعة مليارات سنة ($\num{1.26e17}$~s) يفقد المريخ
$\num{3.8e43}$ ذرّة، $\num{1.9e43}$ جزيئًا من الماء، $\qty{5.7e17}{kg}$؛
على مساحة $4\pi(\num{3.3895e6})^2 = \qty{1.44e14}{m^2}$، أي طبقةً
عمقها \qty{3.9}{m}. (ب)~عشرين ضعفًا أكثر، \qty{78}{m}. (ج)~حتى
معدّل الميل المرتفع لا يبلغ سوى الطرف الأدنى من المدى 100 إلى
\qty{1500}{m}؛ وذهب الباقي إلى القشرة والجليد، أو فُقد
أسرع حين كانت الشمس فتيّة.
\end{solution}

\begin{solution}{exo:b3:planetary-atmospheres:8}
(أ)~$m = pA/g = 636\times\num{1.44e14}/3.71 = \qty{2.5e16}{kg}$.
(ب)~$0.1\times\num{1.26e17} = \qty{1.3e16}{kg}$، أي نصف
الغلاف الجوي الحالي؛ وكلّه في $\num{2.5e16}/0.1 = \qty{2.5e17}{s}$،
\num{7.8}~مليار~سنة. (ج)~كانت الشمس الفتيّة أكثر نشاطًا وعواصفها
أكثر تكرارًا، ويزداد الفقد خلال العواصف؛ قنوات أخرى
(الهيدروجين، والتذرية، والإفلات الكيميائي الضوئي) لم تُحسب.
\end{solution}

\begin{solution}{exo:b3:planetary-atmospheres:9}
(أ)~$S(1 - 0.193)/4 = 291$ يعطي $S = \qty{1442}{W.m^{-2}}$، أي 1.06
ضعف قيمة الأرض. (ب)~كانت الزهرة اليوم ستمتصّ $0.807\times2601/4 =
\qty{525}{W.m^{-2}}$؛ والزهرة الفتيّة تتلقّى $0.7\times2601 = \qty{1821}{W.m^{-2}}$
وتمتصّ \qty{367}{W.m^{-2}}: كلتاهما فوق الحدّ، لذا
كانت الزهرة الفتيّة الخالية من السحب ستدخل \omterm{def:b3:planetary-atmospheres:runaway}{الدفيئة الجامحة} منذ البداية. (ج)~ترفع السحب
الألبيدو؛ وقد وجدت نماذج ثلاثية الأبعاد تضمّ السحب مسارًا
نحو زهرة فتيّة صالحة للحياة.
\end{solution}

\begin{solution}{exo:b3:planetary-atmospheres:10}
(أ)~$H = kT/mg = \num{1.381e-23}\times1000/(16\times\num{1.661e-27}\times9.21)
= \qty{56}{km}$. (ب)~$n_\mathrm{exo} = 1/(\sigma_c H) = 1/(\num{3e-19}\times
\num{5.64e4}) = \qty{5.9e13}{m^{-3}}$; $z = 200 + 56.4\ln(\num{1e16}/
\num{5.9e13}) = 200 + 56.4\times5.13 = \qty{490}{km}$. (ج)~تضعف الجاذبية
كلّما ارتفعنا، فيجعل ارتفاع المقياس يزداد مع الارتفاع، فتتناقص الكثافة
ببطء أكبر ممّا افترض: تتحرّك \omterm{def:b3:planetary-atmospheres:exobase}{قاعدة الغلاف الخارجي} إلى الأعلى.
\end{solution}

\begin{solution}{exo:b3:planetary-atmospheres:11}
(أ)~$\Phi = \num{3.0e13}\times\num{3e-3} = \qty{9e10}{cm^{-2}.s^{-1}} =
\qty{9e14}{m^{-2}.s^{-1}}$؛ على مساحة $4\pi(\num{6.371e6})^2 =
\qty{5.10e14}{m^2}$، $\num{4.6e29}$ جزيئًا من $\ce{H2}$ في الثانية.
(ب)~يحتوي المحيط على $\num{1.4e24}/18.015\times\num{6.022e23} =
\num{4.7e46}$ جزيئًا من الماء، يقابل كلّ منها جزيء $\ce{H2}$ واحد: يدوم
$\qty{1.0e17}{s} = \num{3.2}\text{ مليار سنة}$، قريبة من عمر
الأرض، ولهذا تُحدّد نسبة الخلط \num{3e-3} حدّ
\omterm{def:b3:planetary-atmospheres:runaway}{الدفيئة الرطبة}.
\end{solution}

\begin{solution}{exo:b3:planetary-atmospheres:12}
(أ)~$p_1 = g/D\kappa = 9.81/(1.66\times0.005) = \qty{1180}{Pa}$.
(ب)~$1/T = 1/273.16 - \ln(1182/611.657)\times461.5/\num{2.5e6}$، $T =
\qty{282.5}{K}$، و$\sigma T^4 = \qty{361}{W.m^{-2}}$. (ج)~الحدّ الحقيقي
أدنى، لذا يكون مستوى الإشعاع أبرد وعند ضغط أقلّ: بخار الماء
الحقيقي أكثر عتامةً من هذا الغاز الرمادي، وأنّ $\kappa$ الفعّال له
أكبر.
\end{solution}

\begin{solution}{pb:b3:planetary-atmospheres:1}
\textbf{1.} $g = GM/R^2 = \num{6.674e-11}\times\num{5.2e24}/(\num{6.2e6})^2 =
\qty{9.03}{m.s^{-2}}$; $\bar\rho = M/\tfrac43\pi R^3 = \qty{5210}{kg.m^{-3}}$.

\textbf{2.} $v_\mathrm{esc} = \sqrt{2GM/r} = \sqrt{2\times\num{3.47e14}/\num{6.7e6}}
= \qty{10.18}{km.s^{-1}}$.

\textbf{3.} $u = \sqrt{2kT/m}$: الهيدروجين \qty{4.81}{km.s^{-1}}، الهيليوم
\qty{2.41}{km.s^{-1}}، الأكسجين \qty{1.21}{km.s^{-1}}.

\textbf{4.} $\lambda = (v_\mathrm{esc}/u)^2$: 4.49, 17.8, 71.2.

\textbf{5.} لا يفلت إلا الهيدروجين: كسره المفلت هو 0.06، وكسر الهيليوم
\num{3.5e-7}، وكسر الأكسجين \num{9e-30}. لا شيء منها قريب من النظام الهيدروديناميكي،
إذ $\lambda \approx 1.5$ إلى 3.

\textbf{6.} $\lambda = 4.49\times1400/3000 = 2.09$: ضمن الانتقال؛
كان الهيدروجين سيندفع بعيدًا هيدروديناميكيًا، وربما جرّ معه
الغازات الأثقل.

\textbf{7.} $(1 + 4.49)\,e^{-4.49} = 0.0618$.

\textbf{8.} $n u/2\sqrt\pi = \num{2e11}\times4806/3.545 = \qty{2.71e14}{m^{-2}.s^{-1}}$;
$\Phi_\mathrm J = \num{2.71e14}\times0.0618 = \qty{1.68e13}{m^{-2}.s^{-1}}$.

\textbf{9.} $4\pi(\num{6.7e6})^2 = \qty{5.64e14}{m^2}$; $\num{9.5e27}$ ذرّة في
الثانية.

\textbf{10.} $\num{1.4e24}/18.015 = \qty{7.8e22}{mol}$، $\num{4.68e46}$ جزيئًا،
$\num{9.36e46}$ ذرّة هيدروجين.

\textbf{11.} $\num{9.36e46}/\num{9.5e27} = \qty{9.9e18}{s}$، أي نحو 310 مليارات
سنة.

\textbf{12.} $\lambda_\mathrm D = 4.49\times2.014/1.008 = 8.96$؛ الكسر المفلت
$9.96\,e^{-8.96} = \num{1.28e-3}$؛ الكفاءة النسبية $0.707\times
\num{1.28e-3}/0.0618 = 0.015$.

\textbf{13.} $H_a = kT/m_a g = \num{1.381e-23}\times200/(44.01\times\num{1.661e-27}
\times9.03) = \qty{4.18}{km}$.

\textbf{14.} $b = \qty{1.2e21}{m^{-1}.s^{-1}}$، لذا
\[
\Phi_\mathrm{DL} = \frac{b f(m_a - m_i)}{m_a H_a}
= \frac{\num{1.2e21}\times\num{1e-3}\times(42.0/44.0)}{4185}
= \qty{2.74e14}{m^{-2}.s^{-1}}.
\]

\textbf{15.} $\num{3.0e13}\times\num{1e-3}\times9.03/9.81 =
\qty{2.76e10}{cm^{-2}.s^{-1}} = \qty{2.76e14}{m^{-2}.s^{-1}}$: المقدار نفسه
في حدود 1\,\%.

\textbf{16.} $4\pi(\num{6.2e6})^2 = \qty{4.83e14}{m^2}$; $\num{1.32e29}$ جزيئًا
من $\ce{H2}$، $\num{2.64e29}$ ذرّة في الثانية.

\textbf{17.} يوفّر حدّ الانتشار 28 ضعفًا مما تسمح به \omterm{def:b3:planetary-atmospheres:exobase}{قاعدة الغلاف الخارجي} عند
كثافتها الحالية. يتراكم الهيدروجين في الغلاف الحراري حتى
يتساوى دفق جينز مع الإمداد: على المدى الطويل يكون \omterm{def:b3:planetary-atmospheres:hydrodynamic}{الإفلات
المحدود بالانتشار}، كما على الأرض.

\textbf{18.} $\num{2e11}\times28 = \qty{5.6e12}{m^{-3}}$.

\textbf{19.} $\num{4.68e46}/\num{1.32e29} = \qty{3.54e17}{s}$، أي 11.2 مليار سنة.

\textbf{20.} $S = 0.30/0.60^2 = 0.833$.

\textbf{21.} $d = \sqrt{L/S_\mathrm{eff}}$: الحافة الداخلية $\sqrt{0.30/0.939} =
\qty{0.565}{au}$، الحافة الخارجية $\sqrt{0.30/0.290} = \qty{1.02}{au}$. كيبوس،
عند \qty{0.60}{au}، يقع داخل النطاق، قريبًا من حافته الداخلية.

\textbf{22.} $0.939/0.833 = 1.13$: زيادة في السطوع بمقدار 13\,\%.

\textbf{23.} $0.833\times1361\times0.7/4 = \qty{198}{W.m^{-2}}$, $T_\mathrm{eq} =
(198/\sigma)^{1/4} = \qty{243}{K}$؛ أدنى بكثير من
حدّ \qty{291}{W.m^{-2}} (محسوب لنجم شبيه بالشمس، لكنّ الهامش
واسع).

\textbf{24.} أسرع بثلاث مرات: 3.7 مليار سنة.

\textbf{25.} \textbf{\mbox{\textdir TLT 11}~مليار~سنة}.
\end{solution}
